\label{gu:cov:section}
To pass from the Shalika period identity for big theta lifts to
arbitrary irreducible admissible representations of $\GU(W)$,
we prove that each such representation is a quotient of a suitable
big theta lift after varying the splitting character. The proof
combines analytic equal-rank theta dichotomy, highly ramified
stability, and a comparison of ordinary and similitude theta
correspondence.

\subsection{Analytic equal-rank theta and stability}\label{ana:section}

\subsubsection{Forms, characters, and normalization}

We use the extension $E/F$ and its conjugation $x\mapsto\bar x$
fixed in the introduction. The associated quadratic character of
$F^\times$ is denoted by $\omega_{E/F}$. We use the fixed nontrivial
additive character $\psi$ of $F$ and choose $\delta\in E^\times$ with
$\bar\delta=-\delta$.

For $n=2,4$, let $W_n$ be a hyperbolic skew-Hermitian space of dimension
$n$ over $E$. Let $V_n^+$ be a hyperbolic Hermitian space of dimension
$n$, and let $V_n^-$ be the other Hermitian space of that dimension.
For either type of form set
\[
 \operatorname{disc}(A)=(-1)^{n(n-1)/2}\det(A),
\]
where the determinant is taken modulo $N_{E/F}(E^\times)$. Thus
$\operatorname{disc}(V)\in F^\times/N_{E/F}(E^\times)$ in the
Hermitian case and
$\operatorname{disc}(W)\in\delta^n F^\times/N_{E/F}(E^\times)$
in the skew-Hermitian case; see \cite[\S2.1]{GanIchino}. Define
\[
 \epsilon(V)=\omega_{E/F}(\operatorname{disc}(V)),\qquad
 \epsilon(W)=\omega_{E/F}(\delta^{-n}\operatorname{disc}(W)).
\]
Thus $\epsilon(W_n)=\epsilon(V_n^+)=1$ and
$\epsilon(V_n^-)=-1$. We also put $W_0=V_0^+=0$ and interpret their
unitary groups as the trivial group. There is no space $V_0^-$.

The tensor form is fixed in \eqref{gu:eq:tensor-symplectic}.
Gan--Ichino use the full-trace tensor form
\cite[\S4.1]{GanIchino}. The change of central
coordinate $(z,t)\mapsto(z,-t/2)$ identifies their convention with ours.
Accordingly, all their analytic factors in this section are evaluated
using the fixed character
\begin{equation}\label{ana:psi}
 \psi_0(t)=\psi(-t/2).
\end{equation}

Let $\chi:E^\times\to\C^\times$ be a smooth unitary character satisfying
$\chi|_{F^\times}=1$. The splitting pair $(\chi,1)$ is permissible
for all the even-dimensional spaces just specified. For the full-trace
form with $(\chi,1,\psi_0)$, equivalently our form with
$(\chi,1,\psi)$, write $\omega_{V,W}$ for the ordinary Weil
representation of $\U(V)\times\U(W)$, and put
\begin{equation}\label{ana:big-lift}
 T_{V,W}(\pi)
    =(\omega_{V,W}\otimes\pi^\vee)_{\U(W)}.
\end{equation}
We write $T^\chi_{V,W}$ when the dependence on $\chi$ is needed.
The subscript in \eqref{ana:big-lift} denotes smooth coinvariants.

\subsubsection{Conservation and the analytic sign}

\begin{theorem}[Equal-rank dichotomy and finite length]
\label{ana:dichotomy}
Let $n=2$ or $4$, and let $\pi\in\Irr(\U(W_n))$.
Exactly one of
\[
 T_{V_n^+,W_n}(\pi),\qquad T_{V_n^-,W_n}(\pi)
\]
is nonzero. Both big theta lifts have finite length as smooth
representations of their respective target groups.
\end{theorem}

\begin{proof}
Fix the splitting character $\chi$ in the two compatible even
Hermitian Witt towers. Let $m^+(\pi)$ and $m^-(\pi)$ denote the
dimensions of first occurrence. Sun--Zhu's conservation relation gives
\[
 m^+(\pi)+m^-(\pi)=2n+2;
\]
see \cite[Theorem~1.10 and Proposition~1.6]{SunZhu}.
Both indices are even. If both were at most $n$, their sum would be
at most $2n$. If both exceeded $n$, their sum would be at least
$2n+4$. Hence precisely one is at most $n$. Kudla's persistence
principle, recalled in \cite[\S1.5]{SunZhu}, now proves the assertion
about occurrence. 
Finite length here is the finiteness theorem for the ordinary
\emph{big} theta lift due to Kudla, as recalled in
\cite[Introduction]{GanTakeda} and \cite[\S4.2]{GanIchino}.
The finiteness theorem holds in all residue characteristics.
Consequently a nonzero big theta lift has an irreducible quotient. \end{proof}

Write
$\gamma_{\mathrm{LR}}(s,\pi,\eta,\psi_0)$ for the analytic doubling
factor, for a smooth character $\eta$ of $E^\times$, in the normalization
of \cite[\S10.1 and Appendix~A.5]{GanIchinoFormal}. It is defined by the
normalized local functional equation. 
\begin{proposition}[Tempered analytic criterion]\label{ana:criterion}
Let $n=2$ or $4$ and let $\pi$ be an irreducible tempered
representation of $\U(W_n)$. The analytic factor
$\gamma_{\mathrm{LR}}(s,\pi,\chi^{-1},\psi_0)$ is finite and nonzero
at $s=\tfrac12$. For $V=V_n^+$ or $V_n^-$,
\begin{equation}\label{ana:sign}
 T_{V,W_n}(\pi)\ne0
 \quad\Longleftrightarrow\quad
 \gamma_{\mathrm{LR}}\left(\frac12,\pi,\chi^{-1},\psi_0\right)
       =\omega_\pi(-1)\epsilon(V).
\end{equation}
\end{proposition}

\begin{proof}
The proof uses the analytic argument of \cite[Appendix~A.5]{GanIchinoFormal}.
Let $\mathbf W=W_n\oplus W_n^{-}$, where $W_n^{-}$ has the negative
of the form on $W_n$, and let $I(0,\chi)$ be the normalized Siegel
principal series of $\U(\mathbf W)$ at the center.
The two Siegel--Weil submodules satisfy
\begin{equation}\label{ana:siegel-weil-sum}
 I(0,\chi)=R(V_n^+)\oplus R(V_n^-).
\end{equation}
This is \cite[Proposition~4.1(ii), p.~961]{HKS96}, transported to
the conventions of \cite{GanIchinoFormal}.

Let $\beta_W$ be the nonsingular auxiliary map used to normalize the
doubling intertwiner. By \cite[Appendix~A.4, Lemma~A.3]{GanIchinoFormal},
$M^{\mathrm{LR}}_{\psi_0}(0,\chi)$ acts on $R(V)$ by
\begin{equation}\label{ana:intertwiner-eigenvalue}
 \lambda_V
  =\chi(\det\beta_W)\chi(\delta)^n
        \epsilon(V)\epsilon(W_n)^n.
\end{equation}
The change from the Harris--Kudla--Sweet normalization is computed
in that appendix; the original scalar calculation is
\cite[Proposition~6.8, p.~974]{HKS96}.

For a matrix coefficient $f$ of $\pi^\vee$, the first functional
equation in \cite[Appendix~A.5]{GanIchinoFormal} is
\begin{align}
 &Z\bigl(0,M^{\mathrm{LR}}_{\psi_0}(0,\chi)\mathcal F,f\bigr)
 \notag\\
 &\quad=
   \epsilon(W_n)^{n+1}\omega_\pi(-1)\chi(\det\beta_W)
   \gamma_{\mathrm{LR}}\left(\frac12,\pi,\chi^c,\psi_0\right)
   Z(0,\mathcal F,f),
 \label{ana:functional-equation}
\end{align}
where $\chi^c(x)=\chi(\bar x)$.
For tempered $\pi$, the integrals at the center converge absolutely;
the relevant convergence statement is
\cite[Lemma~9.5(ii)]{GanIchinoFormal}, used expressly in Appendix~A.5.
The resulting family of central zeta functionals is nonzero. Indeed,
one can take a section supported on a sufficiently small compact
neighborhood in the open doubling orbit, and a matrix coefficient
nonzero at the identity.

Finiteness of the central gamma value follows
directly from the meromorphic functional equation. Extend such a
section to a flat holomorphic section near the center. Its zeta
integral is nonzero there, and the normalized intertwiner is regular
at the center by \cite[Proposition~8.1 and Appendix~A]{GanIchinoFormal}.
The functional equation therefore makes the gamma factor holomorphic
at $\tfrac12$. Its value cannot be zero: the operator
$M^{\mathrm{LR}}_{\psi_0}(0,\chi)$ is invertible, as is also evident
from \eqref{ana:siegel-weil-sum} and its two nonzero eigenvalues
\eqref{ana:intertwiner-eigenvalue}. A zero gamma value in
\eqref{ana:functional-equation} would consequently make every
central zeta integral vanish.

By \eqref{ana:siegel-weil-sum}, the zeta functional is nonzero on
at least one $R(V)$. Its restriction to that summand yields theta
occurrence by the doubling see-saw
\cite[Proposition~3.1, p.~960]{HKS96}. The existence hypothesis
on an irreducible theta quotient in that proposition is supplied by
the big-lift finite length in Theorem~\ref{ana:dichotomy}.
Thus this $V$ is the unique nonzero tower in that theorem.
On this summand, \eqref{ana:functional-equation} forces equality
between \eqref{ana:intertwiner-eigenvalue} and the scalar on the
right side of the functional equation. The two eigenvalues have
opposite signs, so this equality also singles out that summand.

Finally, $\chi^c=\chi^{-1}$ because $x\bar x\in F^\times$,
and $\chi(\delta)^n=1$ because $\delta^2\in F^\times$ and $n$ is
even. We have $\epsilon(W_n)=1$ and $\omega_\pi(-1)^2=1$.
Canceling the common factor $\chi(\det\beta_W)$ gives
\eqref{ana:sign}.
\end{proof}

\subsubsection{Highly ramified comparison}

Let $\mathcal O_E$ be the ring of integers of $E$ and $\mathfrak p_E$
its maximal ideal. Put $U_E^0=\mathcal O_E^\times$ and
$U_E^A=1+\mathfrak p_E^A$ for $A\ge1$. For a smooth character
$\eta$ of $E^\times$, write
\[
 a(\eta)=\min\{A\ge0:\eta|_{U_E^A}=1\}.
\]
In particular, $a(\eta^{-1})=a(\eta)$.

\begin{proposition}[Analytic stability with a fixed central character]
\label{ana:stability}
Let $n=2$ or $4$, and let $\pi_1,\pi_2\in\Irr(\U(W_n))$
have the same central character on $E^1$.
For the fixed additive character $\psi_0$, there is an integer
$N=N(\pi_1,\pi_2,\psi_0)$ such that
\begin{equation}\label{ana:stable-gamma}
 \gamma_{\mathrm{LR}}(s,\pi_1,\eta,\psi_0)
   =\gamma_{\mathrm{LR}}(s,\pi_2,\eta,\psi_0)
\end{equation}
as meromorphic functions of $s$, for every smooth character
$\eta$ of $E^\times$ satisfying $a(\eta)>N$.
\end{proposition}

\begin{proof}
This is the unitary case of \cite[Proposition~5.3]{FinisLapid}.
That proposition concerns the analytic Piatetski-Shapiro--Rallis
factors explicated by Lapid--Rallis, and requires equality of central
characters in the unitary case. The group, additive character, and
normalization are fixed in the comparison. The unitary normalization
correction $\epsilon(W_n)^{n+1}$ used in \cite[\S10.1]{GanIchinoFormal}
equals one for the present source; the other common normalizing
factors cancel~\cite[Remarks~5.3--5.5]{Kakuhama19}. Thus the comparison is precisely
\eqref{ana:stable-gamma} in the notation above.
The proposition applies to all sufficiently ramified characters of
$E^\times$, so in particular it applies after imposing
$\eta|_{F^\times}=1$.
\end{proof}

\begin{lemma}[Unbounded conductors]\label{ana:characters}
The smooth unitary characters of $E^\times/F^\times$ have
unbounded conductors.
\end{lemma}

\begin{proof}
The group $Q=E^\times/F^\times$ is compact and totally
disconnected. Indeed, valuations are taken modulo
$e(E/F)\mathbf Z$, and the remaining quotient of unit groups is
compact. Hilbert's Theorem~90 identifies $Q$ with the infinite
group $E^1$ through $x\mapsto x/\bar x$.
Its finite-quotient characters therefore separate points and
form an infinite set. They are unitary, since they have finite
image. For each fixed $A$, a character of conductor at most $A$
factors through the finite group
\[
 E^\times/(F^\times U_E^A).
\]
There are only finitely many such characters. Their conductors
cannot be bounded.
\end{proof}

\subsection{Inducing blocks and the principal-series base case}
\label{blk:section}

\subsubsection{Normalized Jacquet modules}

We retain the splitting characters $(\chi,1)$ and the additive character
fixed above.  Write $G_n=\U(W_n)$ and $H_m=\U(V_m)$, where $W_n$ is
hyperbolic skew-Hermitian and $V_m$ is Hermitian.  In this section all
dimensions are even.  We abbreviate the Weil representation to
$\omega_{V_m,W_n}$ and use $T_{V_m,W_n}(\pi)$ for its ordinary big theta
lift.  All groups denote their groups of $F$-points.  We use normalized
parabolic induction and normalized Jacquet modules:
\[
 i_P^G(\rho)=\operatorname{Ind}_P^G
       (\delta_P^{1/2}\rho),
 \qquad
 r_P(A)=\delta_P^{-1/2}\otimes A_{N_P}.
\]
Thus normalized Frobenius reciprocity uses $r_P$, whereas an induced
representation in the first argument of $\Hom$ uses the opposite
Jacquet module:
\begin{align}
 \Hom_G(A,i_P^G B)&\simeq\Hom_M(r_P A,B),
       \label{blk:first-adjunction}\\
 \Hom_G(i_P^G A,B)&\simeq\Hom_M(A,r_{\overline P}B).
       \label{blk:second-adjunction}
\end{align}
Here $A$ may be smooth and $B$ admissible in
\eqref{blk:second-adjunction}.  For these hypotheses the latter identity
follows from smooth duality, ordinary Frobenius reciprocity, and
Casselman's Jacquet pairing:
\[
\begin{aligned}
 \Hom_G(i_P A,B)
 &\simeq\Hom_G(B^\vee,i_P(A^\vee))\\
 &\simeq\Hom_M(r_P(B^\vee),A^\vee)\\
 &\simeq\Hom_M\bigl(A,(r_P(B^\vee))^\vee\bigr)
  \simeq\Hom_M(A,r_{\overline P}B).
\end{aligned}
\]
The duality for induction applies to smooth $A$ because $P\backslash G$
is compact, and the admissible modules in this calculation are reflexive
\cite[Proposition~3.1.2, Theorems~3.2.4 and~3.3.1,
Corollary~4.2.5]{Casselman}.
For representations of a product group, these identities apply with
admissibility for the product representation.

Let $P_k$ be the parabolic of $G_n$ stabilizing a $k$-dimensional
isotropic subspace $X_k$, with Levi
\[
 L_k=\GL(X_k)\times G_{n-2k}.
\]
If $Y_a\subset V_m$ is isotropic of dimension $a$, let $Q'_a$ be its
stabilizer in $H_m$.  Write
$X_k=X_{k-a}\oplus X'_a$, and denote by $Q_{k-a,a}$ the parabolic of
$\GL(X_k)$ stabilizing $X_{k-a}$.  A subscript of dimension zero denotes
the trivial group or zero space, as appropriate.

\begin{lemma}[The equal-rank boundary exponents]
\label{blk:kudla}
For $g\in\GL_k(E)$, $c\in\GL_a(E)$, and
$(d,b)\in\GL_{k-a}(E)\times\GL_a(E)$, the relevant modulus characters
are
\begin{align}
 \delta_{P_k}(g)&=\abs{\det g}_E^{n-k},
       \label{blk:source-modulus}\\
 \delta_{Q'_a}(c)&=\abs{\det c}_E^{m-a},
       \label{blk:target-modulus}\\
 \delta_{Q_{k-a,a}}(d,b)
   &=\abs{\det d}_E^a\abs{\det b}_E^{-(k-a)}.
       \label{blk:gl-modulus}
\end{align}
The normalized Jacquet module $r_{P_k}(\omega_{V_m,W_n})$ has a
decreasing filtration
\[
 R^0\supset R^1\supset\cdots\supset R^k\supset R^{k+1}=0,
 \qquad J^a=R^a/R^{a+1}.
\]
If $V_m$ has no isotropic $a$-space, then $J^a=0$.  Otherwise $J^a$ is
normalized induction from
\[
 (Q_{k-a,a}\times G_{n-2k})\times Q'_a
   \ \subset\ L_k\times H_m
\]
of a representation with underlying space
\[
 \cS\bigl(\operatorname{Isom}_{E,c}(X'_a,Y_a)\bigr)
       \otimes\omega_{V_{m-2a},W_{n-2k}}.
\]
Here $\operatorname{Isom}_{E,c}$ denotes conjugate-linear $E$-isomorphisms;
thus $xb$ means composition with $b$ on the source.
On this coefficient space the unmatched factor $\GL(X_{k-a})$ acts
by
\begin{equation}
 \chi(\det d)\abs{\det d}_E^{(m-n+k-a)/2}.
 \label{blk:boundary-character}
\end{equation}
The matched source and target factors act on the first tensor factor by
\[
 ((b,c)f)(x)=\chi(\det b)f(c^{-1}xb),
\]
and the two remaining unitary groups act on the smaller Weil
representation.  In particular, for $m=n$ every boundary term $a<k$
has strictly positive unmatched real exponent $(k-a)/2$.
\end{lemma}

\begin{proof}
The off-diagonal part of the unipotent radical of $P_k$ contributes
$\abs{\det g}_E^{n-2k}$ to its modulus.  The Hermitian or skew-Hermitian
$k$-by-$k$ part contributes $\abs{\det g}_E^k$, giving
\eqref{blk:source-modulus}.  The same calculation gives
\eqref{blk:target-modulus}; the action $x\mapsto dxb^{-1}$ on
$\operatorname{Hom}_E(X'_a,X_{k-a})$ gives
\eqref{blk:gl-modulus}.  These are determinant identities over $F$ and
apply in residue characteristic two as well.

At rank zero, the mixed
Schr\"odinger model has scalar factor
$\chi(\det g)\abs{\det g}_E^{m/2}$ before taking normalized Jacquet
modules.  Dividing by $\delta_{P_k}^{1/2}$ gives
\[
 \chi(\det g)\abs{\det g}_E^{m/2-(n-k)/2}
   =\chi(\det g)\abs{\det g}_E^{(m-n+k)/2}.
\]
The source-space filtration is \cite[Lemma~3.1]{GanTakeda}; compare
\cite[Appendix~C.1, Lemma~C.2]{GanIchinoFormal} with the two unitary spaces
interchanged.  The parameter $\epsilon_0$ there is zero for a unitary
pair.  Normalized induction from these parabolics then gives
\eqref{blk:boundary-character}.  Setting $m=n$ proves the last assertion.
\end{proof}

Only the following boundary terms occur in our argument:
\[
\begin{array}{c|c|c|c}
 n&k&a&\text{unmatched real exponent}\\ \hline
 2&1&0&\tfrac12\\
 4&1&0&\tfrac12\\
 4&2&0&1\\
 4&2&1&\tfrac12
\end{array}
\]
We will repeatedly use a simple consequence of the decreasing
filtration.  If $\Hom(J^a,B)=0$ for all $a<k$, then restriction is
injective:
\begin{equation}
 \Hom(R^0,B)\longrightarrow\Hom(J^k,B).
 \label{blk:restriction-injective}
\end{equation}
Indeed, a map vanishing on $J^k=R^k$ factors through $R^0/R^k$, and its
successive restrictions vanish along the finite filtration of that
quotient.  In particular, if $J^k=0$, the first Hom space is zero.

\subsubsection{The zero-exponent principal series}

\begin{proposition}
\label{blk:principal-split}
\label{blk:zero}
Let $n=2\ell\in\{2,4\}$, and let $\alpha_1,\ldots,\alpha_\ell$ be
smooth unitary characters of $E^\times$.  Every irreducible constituent
$\pi$ of
\[
 I=i_B^{G_n}(\alpha_1\boxtimes\cdots\boxtimes\alpha_\ell)
\]
satisfies
\[
 T_{V_n^-,W_n}(\pi)=0,
 \qquad T_{V_n^+,W_n}(\pi)\ne0
\]
for every permissible unitary $\chi$.
\end{proposition}

\begin{proof}
The representation $I$ is admissible, of finite length, and unitary
\cite[Proposition~2.3 and Theorem~2.8]{BZ77}
\cite[Proposition~3.1.4]{Casselman}.  It is therefore semisimple.
Here is the elementary argument relevant to this use.  If $A\subset I$
is invariant and $v\in I^K$ for a compact open subgroup $K$, decompose
$v=a+b$ orthogonally in the finite-dimensional space $I^K$, with
$a\in A^K$ and $b\perp A^K$.  Averaging over $K$ shows that
\[
 \langle b,u\rangle=\langle b,e_Ku\rangle=0
       \qquad(u\in A).
\]
Thus $I=A\oplus A^\perp$.  Finite length now gives a finite direct sum
of irreducibles; see also
\cite[Propositions~2.1.5 and~2.1.14]{Casselman}.  In particular,
$\pi\hookrightarrow I$, even when $I$ is reducible.

Group the characters in the source Siegel Levi.  Normalized
transitivity gives
\begin{equation}
 \rho=i_{B_\ell}^{\GL_\ell(E)}
       (\alpha_1\boxtimes\cdots\boxtimes\alpha_\ell),
 \qquad I=i_{P_\ell}^{G_n}(\rho).
 \label{blk:grouped-principal}
\end{equation}
The center of $\GL_\ell(E)$ acts unitarily on $\rho$.  For $\ell=2$
the geometric lemma gives
\begin{equation}
 r_{\overline B_2}(\rho)^{\mathrm{ss}}
  =(\alpha_1\boxtimes\alpha_2)
       \oplus(\alpha_2\boxtimes\alpha_1).
 \label{blk:opposite-principal}
\end{equation}
This statement about composition factors applies when the characters
coincide as well \cite[Lemma~2.12]{BZ77}.

Suppose $T_{V_n^-,W_n}(\pi)\ne0$.  By
Theorem~\ref{ana:dichotomy} it has finite length, so it has an
irreducible quotient $\sigma$. Composing the corresponding nonzero
map $\omega_{V_n^-,W_n}\to\pi\boxtimes\sigma$ with the embedding
of $\pi$ in \eqref{blk:grouped-principal} gives, by
\eqref{blk:first-adjunction},
\begin{equation}
 \Hom_{\GL_\ell(E)\times H_n}
   \bigl(r_{P_\ell}\omega_{V_n^-,W_n},
               \rho\boxtimes\sigma\bigr)\ne0.
 \label{blk:principal-hom}
\end{equation}

If $\ell=1$, the rank-zero source character is
$\chi\abs{\cdot}_E^{1/2}$, whereas $\rho=\alpha_1$ is unitary.
Consequently the Hom space from $J^0$ is zero.

If $\ell=2$, the source $\GL_2$ acts on $J^0$ by
$\chi(\det)\abs{\det}_E$.  On a scalar $zI_2$ its absolute value is
$\abs{z}_E^2$, whereas the center acts unitarily on $\rho$.
Thus the rank-zero Hom space again vanishes.  For rank one apply
\eqref{blk:second-adjunction} to the source $B_2$ and target $Q'_1$
occurring in Lemma~\ref{blk:kudla}.  The unmatched source $E^\times$
acts in the coefficient by $\chi\abs{\cdot}_E^{1/2}$, while the first
characters of \eqref{blk:opposite-principal} are unitary.  This excludes
the rank-one term.  Possible extensions between the composition factors
in \eqref{blk:opposite-principal} do not affect the argument: a character
with positive real exponent cannot occur in a finite-length module all
of whose characters have real exponent zero.

By \eqref{blk:restriction-injective}, the nonzero Hom space in
\eqref{blk:principal-hom} must be nonzero on $J^\ell$.  But $V_n^-$
has Witt index $\ell-1$, so $J^\ell=0$.  This contradiction proves
vanishing in the nonsplit tower.  Equal-rank dichotomy
(Theorem~\ref{ana:dichotomy}) then gives nonvanishing in the split tower.
\end{proof}

\subsubsection{Positive inducing exponents}

\begin{lemma}[Removal of inducing blocks in dimensions two and four]
\label{blk:remove}
Let $n\in\{2,4\}$.  Suppose that an irreducible representation $\pi$ of
$G_n$ embeds in
\begin{equation}
 i_P^{G_n}\bigl(
   \tau_1\abs{\det}_E^{-s_1}\boxtimes\cdots\boxtimes
   \tau_r\abs{\det}_E^{-s_r}\boxtimes\pi_0\bigr),
 \label{blk:anti-standard}
\end{equation}
where $s_1\ge\cdots\ge s_r>0$, each $\tau_i$ is a unitary discrete
series representation of $\GL_{k_i}(E)$, and $\pi_0$ is irreducible
tempered on $G_{n_0}$, with $n_0=n-2\sum_i k_i$.  If
$T_{V_n,W_n}(\pi)\ne0$, then
\[
 V_n\simeq\mathbb H^{\sum_i k_i}\perp V_{n_0},
 \qquad T_{V_{n_0},W_{n_0}}(\pi_0)\ne0.
\]
Here $\mathbb H$ is a Hermitian hyperbolic plane, and
$T_{0,0}(\mathbf1)=\C$.
\end{lemma}

\begin{proof}
If $r=0$, the assertion is immediate.  Put $k=\sum_i k_i>0$ and group
the general linear blocks into a representation $\rho$ of $\GL_k(E)$.
Normalized transitivity turns \eqref{blk:anti-standard} into an
embedding
\[
 \pi\hookrightarrow i_{P_k}^{G_n}(\rho\boxtimes\pi_0).
\]
Choose an irreducible quotient $\sigma$ of the nonzero big lift.
Composing as in the proof of Proposition~\ref{blk:principal-split}
gives a nonzero map
\begin{equation}
 r_{P_k}\omega_{V_n,W_n}\longrightarrow
         \rho\boxtimes\pi_0\boxtimes\sigma.
 \label{blk:removal-hom}
\end{equation}
We first show that its restriction to $J^k$ is nonzero.

If $k=1$, then $\rho=\alpha\abs{\cdot}_E^{-s}$ with $s>0$ and
$\alpha$ unitary.  The only boundary term has source character
$\chi\abs{\cdot}_E^{1/2}$.  Its Hom space to the target of
\eqref{blk:removal-hom} is zero because the real exponents differ.

If $k=2$, necessarily $n=4$ and $n_0=0$.  There are two cases.  First,
suppose
\[
 \rho=\tau\abs{\det}_E^{-s},\qquad s>0,
\]
with $\tau$ unitary discrete series on $\GL_2(E)$.  The center acts on
$\rho$ with absolute value $\abs{z}_E^{-2s}$ at $zI_2$, whereas it acts
on $J^0$ with absolute value $\abs{z}_E^2$.  This excludes rank zero.
For rank one, second adjointness uses $r_{\overline B_2}(\tau)$.
The real exponent on its first $E^\times$ factor is negative whenever
this Jacquet module is nonzero, by Casselman's square-integrability
criterion; this is precisely the opposite-parabolic use of the
criterion in \cite[Appendix~C.1, equations~(C.1)--(C.2)]{GanIchinoFormal}.
After twisting by $\abs{\det}_E^{-s}$ the exponent remains negative,
so it cannot equal the boundary exponent $1/2$.  For example, a
unitary Steinberg twist has
\[
 r_{\overline B_2}\bigl(\operatorname{St}_2\otimes(\alpha\circ\det)\bigr)
 =\alpha\abs{\cdot}_E^{-1/2}
       \boxtimes\alpha\abs{\cdot}_E^{1/2},
\]
which gives the unmatched exponent $-s-1/2$.

Second, suppose
\[
 \rho=i_{B_2}^{\GL_2(E)}
       \bigl(\alpha_1\abs{\cdot}_E^{-s_1}
             \boxtimes\alpha_2\abs{\cdot}_E^{-s_2}\bigr).
\]
The center of $\rho$ has absolute value
$\abs{z}_E^{-(s_1+s_2)}$, excluding $J^0$.  Its opposite Jacquet
composition factors are the two ordered inducing characters and their
permutation.  Thus the unmatched first exponent is $-s_1$ or $-s_2$,
both negative, excluding $J^1$.  This argument uses the entire induced
representation $\rho$ and remains valid when $\rho$ is reducible.

These cases exhaust $k\le2$.  Equation~\eqref{blk:restriction-injective}
now shows that \eqref{blk:removal-hom} is nonzero on $J^k$.  Consequently
$V_n$ contains an isotropic $k$-space and has a decomposition
\[
 V_n\simeq\mathbb H^k\perp V_{n_0}.
\]
If $n_0=0$, this proves the assertion, including the stated nonvanishing
for the zero-dimensional pair.

It remains to consider $n_0=2$.  Then $n=4$ and $k=1$.  The coefficient
of the bottom term is
\[
 \mathcal A^1=
 \cS\bigl(\operatorname{Isom}_{E,c}(X_1,Y_1)\bigr)
       \otimes\omega_{V_2,W_2},
\]
with the actions of Lemma~\ref{blk:kudla}, and $J^1$ is its normalized
induction in the target from $Q'_1$.  Applying
\eqref{blk:second-adjunction} in that target gives a nonzero map
\begin{equation}
 \mathcal A^1\longrightarrow
   (\alpha\abs{\cdot}_E^{-s})\boxtimes\pi_0
                   \boxtimes r_{\overline Q'_1}(\sigma).
 \label{blk:bottom-map}
\end{equation}
This is a map of representations of
\[
 E^\times\times G_2\times E^\times\times\U(V_2).
\]
The Jacquet module in \eqref{blk:bottom-map} is admissible of finite
length.  Choose its composition series, with factors
$\nu_j\boxtimes\sigma_{0,j}$.  Tensoring this series with the fixed
irreducible $(\alpha\abs{\cdot}_E^{-s})\boxtimes\pi_0$ gives a
composition series of the target of \eqref{blk:bottom-map}; the factors
are irreducible by the tensor-product theorem for representations of a
product group \cite[\S1.6]{BZ77}.

Let $C$ be the nonzero image of \eqref{blk:bottom-map}, and choose a
simple quotient $C\twoheadrightarrow S$.  By Jordan--H\"older,
\[
 S\simeq
 (\alpha\abs{\cdot}_E^{-s})\boxtimes\pi_0
                      \boxtimes\nu\boxtimes\sigma_0
\]
for some constituent $\nu\boxtimes\sigma_0$ of the chosen composition series.  Composing
$\mathcal A^1\twoheadrightarrow C\twoheadrightarrow S$ gives a
nonzero map with this irreducible target.  Choose
$f\in\cS(\operatorname{Isom}_{E,c}(X_1,Y_1))$ on which the resulting map is
nonzero.  Fixing $f$ and identifying the two one-dimensional character
factors with $\C$ yields a nonzero
$G_2\times\U(V_2)$-equivariant map
\[
 \omega_{V_2,W_2}\longrightarrow\pi_0\boxtimes\sigma_0.
\]
Therefore $T_{V_2,W_2}(\pi_0)\ne0$, as required.
\end{proof}

\subsubsection{An irreducible principal series with prescribed center}

\begin{lemma}
\label{blk:prescribed-center}
\label{blk:principal}
For every smooth unitary character $\zeta$ of $E^1$ and every
$n\in\{2,4\}$, there exists an irreducible unitary tempered principal
series $\pi^*$ of $G_n$ with central character $\zeta$.  For every
permissible unitary $\chi$,
\[
 T_{V_n^+,W_n}(\pi^*)\ne0.
\]
\end{lemma}

\begin{proof}
Since $E^1\subset\mathcal O_E^\times$, choose $A\ge1$ such that
$\zeta$ is trivial on $E^1\cap(1+\mathfrak p_E^A)$.  The induced
character of the image of $E^1$ in the finite abelian group
$\mathcal O_E^\times/(1+\mathfrak p_E^A)$ extends to that entire group.
Indeed, characters of subgroups of finite abelian groups extend by
successively choosing roots in $\C^\times$ on cyclic generators.
The extension is unitary.  Giving a uniformizer a unitary value then
defines a smooth unitary character $\alpha$ of $E^\times$ extending
$\zeta$.

Put $\ell=n/2$ and initially choose
$(\alpha_1,\ldots,\alpha_\ell)=(\alpha,1,\ldots,1)$. We may twist
the $\alpha_i$ independently by unramified unitary characters $\nu_i$.
These twists are trivial on $E^1$ and preserve
\[
 (\alpha_1\nu_1\cdots\alpha_\ell\nu_\ell)|_{E^1}=\zeta.
\]
They can be chosen so that the resulting character $\beta$ of the
Borel torus is Weyl-regular.  To see this explicitly, write
$\nu_i(\varpi_E)=z_i\in S^1$.  The relative Weyl group acts by signed
permutations, where a sign replaces a character by its inverse
conjugate.  For any nonidentity Weyl element, fixing $\beta$ either is
impossible on $\mathcal O_E^\times$, or imposes a nontrivial equation
\[
 z_i=c z_j,\qquad z_i=c z_j^{-1},\qquad\text{or}\qquad z_i^2=c
\]
with $c\in S^1$ fixed.  Each stabilizer locus is a proper closed subset
of $(S^1)^\ell$ with empty interior.  Their finite union does not cover
the torus.  Choose the $\nu_i$ in its complement.

Let $T$ be the Levi torus of $B$, and let $I=i_B^{G_n}(\beta)$.
The geometric lemma gives one composition
factor $w\beta$ in $r_BI$ for each Weyl element $w$.  Regularity implies
that $\beta$ occurs exactly once.  Hence
\[
 \dim\operatorname{End}_{G_n}(I)
   =\dim\Hom_T(r_BI,\beta)\le1.
\]
The identity endomorphism gives equality.  Since $I$ is unitary,
admissible, and of finite length, it is semisimple by the argument in
Proposition~\ref{blk:principal-split}; the equality forces $I$ to be
irreducible.  This is also
\cite[Theorem~6.6.1]{Casselman}.  Normalized induction preserves
temperedness \cite[Lemma~III.2.3]{Waldspurger03}, so $I$ is tempered.

A scalar $z\in E^1$ corresponds to $(z,\ldots,z)$ in the Borel torus,
and $\delta_B(z)=1$.  Its action on $I$ is consequently
\[
 \prod_{i=1}^{\ell}\alpha_i(z)\nu_i(z)=\zeta(z).
\]
Thus $\pi^*=I$ has the prescribed central character.
Proposition~\ref{blk:principal-split} gives the asserted split
nonvanishing for every permissible $\chi$.
\end{proof}

\subsection{Simultaneous ordinary nonvanishing}
\label{cov:ordinary-section}

Konno's Langlands classification
\cite[Theorem~3.5 and Corollary~3.2]{Konno03}, applied to $\pi^\vee$
and followed by smooth duality, embeds $\pi$ into a normalized
anti-standard induction from an irreducible tempered Levi representation
with strictly negative real exponents on its general linear factors.
For the hyperbolic $\U(W_n)$, $n=2,4$, such a Levi is
\[
 \GL_{k_1}(E)\times\cdots\times\GL_{k_r}(E)
                  \times\U(W_{n_0}),\qquad
 n_0=n-2\sum_i k_i\in\{0,2,4\}.
\]
If a general linear tempered factor is not discrete series, refine it
using the direct-summand statement
\cite[Proposition 2.2]{Konno03}, which is
\cite[Proposition III.4.1]{Waldspurger03}. Exactness and transitivity of
normalized induction preserve the embedding. In these dimensions the
only possible refinement replaces a tempered $\GL_2$ factor by an
induction of two unitary $\GL_1$ characters, both carrying the same
strictly negative twist. Thus we have the embedding required by
Lemma~\ref{blk:remove}, with unitary discrete-series blocks and a
tempered $\pi_0$. The zero-dimensional core is the trivial representation of the trivial
group.

\begin{theorem}[Every sufficiently ramified permissible character works]
\label{cov:ordinary}
Fix the additive character and the hyperbolic source and target spaces.
Let $\pi_1,\ldots,\pi_r$ be a finite family of irreducible admissible
smooth representations of $\U(W_4)$. There is an integer $N$ such that
for every smooth character $\chi:E^\times/F^\times\to\C^\times$ with
$a(\chi)>N$,
\[
 T^{\chi}_{V_4^+,W_4}(\pi_i)\ne0
             \qquad(1\le i\le r).
\]
The assertion also holds for a finite family drawn from both
$\U(W_2)$ and $\U(W_4)$, with each target of the corresponding dimension.
Such characters exist with arbitrarily large conductor.
\end{theorem}

\begin{proof}
Choose a tempered core $\pi_{i,0}$ and an anti-standard embedding as above
for each $\pi_i$. Each positive-dimensional core has a smooth unitary
central character on the compact group $E^1$. By
Lemma~\ref{blk:principal}, choose an irreducible tempered unitary principal
series $\pi_{i,0}^*$ of the \emph{same} unitary group with
\[
 \omega_{\pi_{i,0}^*}=\omega_{\pi_{i,0}}
                  \quad\hbox{on all of }E^1,
\]
whose equal-rank lift is in the hyperbolic tower for every permissible
$\chi$. All these comparators are chosen before imposing a conductor
bound on $\chi$.

For each of these finitely many pairs, the stability theorem
\cite[Proposition~5.3]{FinisLapid} provides a conductor bound with
the same additive character $\psi_0$ and the normalization fixed
in the analytic section. Let $N$ be the maximum of these bounds (and take $N=0$ if there are no positive-dimensional
cores). For every permissible $\chi$ with $a(\chi)>N$, inversion preserves
conductor, so stability applies to $\eta=\chi^{-1}$. The tempered analytic
criterion then gives
\begin{align*}
 \gamma_{\mathrm{LR}}\!\left(\tfrac12,\pi_{i,0},\chi^{-1},
                         \psi_0\right)
 &=\gamma_{\mathrm{LR}}\!\left(\tfrac12,\pi_{i,0}^*,\chi^{-1},
                         \psi_0\right)\\
 &=\omega_{\pi_{i,0}^*}(-1)
  =\omega_{\pi_{i,0}}(-1).
\end{align*}
These central values are finite and nonzero. The same criterion places
each positive-dimensional core in its hyperbolic equal-rank tower.
For a zero-dimensional core the only target is the zero space, and its
lift is nonzero.

For the fixed character $\chi$ under consideration, equal-rank dichotomy
(Theorem~\ref{ana:dichotomy}) provides a nonzero lift of each original
$\pi_i$ in one of the two towers. Lemma~\ref{blk:remove} forces
the target of that lift to be
\[
 \mathbb H^{\sum_j k_{ij}}\perp V_{n_{i,0}},
 \qquad T^{\chi}_{V_{n_{i,0}},W_{n_{i,0}}}(\pi_{i,0})\ne0.
\]
Uniqueness of the core's nonzero equal-rank tower makes $V_{n_{i,0}}$
hyperbolic. Adjoining the hyperbolic summand $\mathbb H^{\sum_j k_{ij}}$ therefore gives
the prescribed $V_4^+$, or $V_2^+$ when the original source has dimension two. In
dimension zero, the target is the sum of the hyperbolic planes
in the decomposition above.
Unbounded conductors of characters of $E^\times/F^\times$ give the last
assertion.
\end{proof}

\subsection{Restriction, reverse theta, and the quotient argument}
\label{fun:gu-passage}

Let $G=G_\theta$, $H=G'_\theta$, $G_1=G_\theta^1$, and
$H_1=(G'_\theta)^1$. The matching group $R$ is as in
Section~\ref{gu:jac:oscillator-section}. We fix permissible splitting data
$D=(\chi_V,\chi_W,\kappa)$ as in Appendix~\ref{gu:app:splitting},
and write $\varpi_D=\omega_{R,D}$ on $\cS=\mathcal S(Z)$. Define
\[
 \Omega_D=\cInd_R^{G\times H}\varpi_D,\qquad
 \Theta_D^{\GU}(\tau)=(\Omega_D\otimes\tau^\vee)_G.
\]
For $D=(\chi,1,1)$, the restriction to $G_1\times H_1$ is the ordinary
Weil representation used above. Ordinary big theta lifts have finite
length by Theorem~\ref{ana:dichotomy}; their irreducible subquotients
are admissible by \cite[Theorem~2.5]{BZ77}. Exactness of compact-open
invariants makes the big lifts admissible as well.

\subsubsection{Restriction and smooth duality}

\begin{lemma}\label{fun:restriction}
For \(\Pi\in\Irr(H)\), there are pairwise inequivalent
irreducible admissible representations \(\pi_i\) of \(H_1\) such that
\[
 \Pi|_{H_1}\simeq\bigoplus_{i=1}^r\pi_i,\qquad
 \Pi^\vee|_{H_1}\simeq\bigoplus_{i=1}^r\pi_i^\vee,\qquad r\in\{1,2\}.
\]
Their isomorphism classes form one \(H\)-conjugacy orbit.
\end{lemma}

\begin{proof}
The scalar center is \(Z=E^\times\). The section
\(d_s=\operatorname{diag}(I_2,sI_2)\) has multiplier \(s\), so
\[
 H_0=ZH_1=\lambda^{-1}(N_{E/F}E^\times),\qquad H/H_0\simeq C_2.
\]
The subgroup \(H_0\) is normal and open. The center acts on \(\Pi\)
by a smooth character \(\omega_\Pi\).

Finite-index Clifford theory makes \(\Pi|_{H_0}\) semisimple of
finite length. \(\Pi\) is generated by one vector
over \(H\), hence is finitely generated over \(H_0\), and has an
irreducible \(H_0\)-quotient \(\rho\). Frobenius reciprocity for the
open subgroup \(H_0\) embeds \(\Pi\) in
\(\operatorname{Ind}_{H_0}^H\rho\). Restricted to \(H_0\), this
induction is the direct sum of the two conjugates of \(\rho\).

Because \(Z\) acts by scalars, the \(H_1\)-submodules and
\(H_0\)-submodules of \(\Pi\) are identical. Irreducibility makes
the restriction constituents one conjugacy orbit. Their common
multiplicity is the dimension of an irreducible projective
representation of an inertia quotient in \(C_2\). Such a
representation is one-dimensional: a matrix representing a
generator has square \(cI\), and rescaling by a square root of
\(c\ne0\) reduces to an ordinary cyclic representation.

The restriction is admissible. If \(J\subset H_1\) is compact open,
choose a compact open \(Z_0\subset Z\) in \(\ker\omega_\Pi\).
The central-product map \(Z\times H_1\to H_0\) is open, and
\(Z_0J\) is compact open in \(H\). Thus
\(\Pi^J=\Pi^{Z_0J}\) is finite-dimensional.

If a functional on \(\Pi\) is \(H_1\)-smooth, a subgroup \(J\) as
above fixes it, and so does \(Z_0J\). It is therefore \(H\)-smooth.
The two smooth subspaces of the algebraic dual coincide, proving
the assertion about contragredients. These are the elementary statements of~\cite[Lemmas~5.2--5.3, \S9]{BGS}.
\end{proof}

\subsubsection{The reverse module}

Define
\[
 M_D(\Pi)=(\cS\otimes\Pi^\vee)_{H_1}.
\]
The \(R\)-action on the tensor product descends across the normal
subgroup \(1\times H_1\), so \(M_D(\Pi)\) carries an action of
\(R/(1\times H_1)\simeq G\). Namely,
\begin{equation}\label{fun:matching-action}
 g[v\otimes\ell]=[\varpi_D(g,h)v\otimes\Pi^\vee(h)\ell],
 \qquad \lambda(h)=\lambda(g).
\end{equation}
The formula is independent of the matching \(h\). The action is
smooth, since the quotient projection has the continuous section
\(g\mapsto(g,d_{\lambda(g)})\).

Lemma~\ref{gu:app:comparison} gives the $G$-equivariant isomorphism
\[
 (\Omega_D\otimes\Pi^\vee)_H\simeq M_D(\Pi),
\]
with quotient map \eqref{gu:asm:quotient-integral}.

\subsubsection{The fixed-data equivalence}

\begin{proposition}\label{fun:fixed-data}\label{gu:asm:assembly}
For \(\Pi\in\Irr(H)\) and fixed \(D\), the following
are equivalent:
\begin{enumerate}
\item Some irreducible admissible \(\tau\) of \(G\) has a quotient
\(\Theta_D^{\GU}(\tau)\twoheadrightarrow\Pi\).
\item \(M_D(\Pi)\ne0\).
\item Some constituent of \(\Pi|_{H_1}\) has nonzero ordinary big
theta lift to \(G_1\) for \(D\).
\end{enumerate}
The reverse module is of finite length and admissible. Any
irreducible quotient of a nonzero \(M_D(\Pi)\) gives the
surjection in (1).
\end{proposition}

\begin{proof}
Lemma~\ref{fun:restriction} and \eqref{fun:matching-action} give
\[
 M_D(\Pi)|_{G_1}
 \simeq\bigoplus_i(\cS\otimes\pi_i^\vee)_{H_1}.
\]
The summands are ordinary big theta lifts. Thus (2) and (3) are
equivalent. Their finite length gives finite length of \(M_D(\Pi)\),
because every strict chain of \(G\)-submodules is a strict chain
of \(G_1\)-submodules.

The reverse module is also admissible. For compact open \(K\subset G\),
\(K\cap G_1\) is compact open in \(G_1\), and
\[
 M_D(\Pi)^K\subset
 \bigl(M_D(\Pi)|_{G_1}\bigr)^{K\cap G_1}.
\]
The right side is finite-dimensional. Compact averaging makes
\(K\)-invariants exact, so every quotient of \(M_D(\Pi)\) is
admissible.

Suppose \(M_D(\Pi)\ne0\), and choose an irreducible quotient
\(p:M_D(\Pi)\twoheadrightarrow\tau\). Lemmas~\ref{fun:tensor}
and \ref{fun:reverse-comparison} give
\[
\begin{aligned}
 \Hom_G(M_D(\Pi),\tau)
 &\simeq\Hom_{G\times H}
       (\Omega_D,\tau\boxtimes\Pi)\\
 &\simeq\Hom_H
       (\Theta_D^{\GU}(\tau),\Pi).
\end{aligned}
\]
The two tensor-evaluation steps use admissibility of \(\Pi\)
and \(\tau\), respectively. The map corresponding to \(p\) is
nonzero and hence surjective. Conversely, a nonzero map in (1)
gives a nonzero element of the left Hom space, proving (2).
\end{proof}

For these groups, some constituent in (3) can be replaced by
every constituent. They form one \(H\)-orbit, and each \(h\in H\)
has a matching \(g\in G\). The operator \(\varpi_D(g,h)\) transports
the associated ordinary theta modules up to conjugation of the
\(G_1\)-action. Their nonvanishing is therefore identical. In
particular, the existence part needs a splitting choice for only
one restriction constituent.

The scalar and multiplier actions for general $D$ are recorded in
Appendix~\ref{gu:app:splitting}, together with the possible multiplier
self-twists.

\subsection{Coverage and Shalika uniqueness}
\label{cov:gu-section}

\begin{theorem}[Finite-family variable-splitting coverage]
\label{cov:gu}\label{gu:asm:coverage}
The spaces $V_4^+$ and $W_4$, the additive character, and the
normalized matching oscillator are fixed throughout. For every finite family
\[
 \Pi_1,\ldots,\Pi_r\in\Irr(\GU(W_4))
\]
there is an integer $N$ such that for every smooth character
$\chi:E^\times/F^\times\to\C^\times$ with $a(\chi)>N$, there are
irreducible admissible smooth representations
$\tau_j\in\Irr(\GU(V_4^+))$ and surjections
\[
 \Theta^{\GU}_{\chi,1,1}(\tau_j)\twoheadrightarrow\Pi_j
                    \qquad(1\le j\le r).
\]
All these characters $\chi$ are unitary. They have unbounded conductor.
In particular, varying only $\chi_V=\chi$ suffices: $\chi_W=1$ and
$\kappa=1$ may remain fixed. 
\end{theorem}

\begin{proof}
Restriction of each $\Pi_j$ to $\U(W_4)$ is a sum of one or two
inequivalent irreducible admissible representations, by the restriction
lemma in the preceding section. Applying Theorem~\ref{cov:ordinary}
to all these constituents yields a common bound $N$ and simultaneous
ordinary nonvanishing for every $\chi$ above it. For each such $\chi$ and each $j$, the fixed-data equivalence
of Proposition~\ref{fun:fixed-data} supplies the required quotient.
Compactness of $E^\times/F^\times$ makes every smooth complex character
of this quotient unitary. 
\end{proof}

\begin{proof}[Proof of Theorem~\ref{thm:main-gu}]
Choose $q:\Theta^{\GU}_{\chi,1,1}(\tau)\twoheadrightarrow\Pi$ by
Theorem~\ref{cov:gu}. Precomposition by a surjective map is injective,
so
\[
 \Hom_{S_B^{\GU}}(\Pi,\sigma\boxtimes\psi_B)
 \lhook\joinrel\longrightarrow
 \Hom_{S_B^{\GU}}(\Theta^{\GU}_{\chi,1,1}(\tau),
                                  \sigma\boxtimes\psi_B).
\]
By Theorem~\ref{gu:thm:hom}, the space on the right is
\[
 \Hom_{H_{\mathrm U}}\!\left(
   \tau^\vee\boxtimes
   \bigl((\sigma^j)^\vee\otimes(\chi^{-1}\circ\det_E)\bigr),
   \mathbf1\right).
\]
Both $\tau^\vee$ and the indicated twist of $(\sigma^j)^\vee$ are irreducible admissible: smooth
contragredients preserve irreducibility in the admissible category,
and character twists preserve it. Theorem~\ref{thm:main-unitary} bounds this Hom space by one.
\end{proof}

